% AUTO-GENERATED by tmp/build_sec34_body.py from the Phase-C paragraph ledgers.
% Source of truth = section{3.1,3.2,3.3,3.4,4.1}_paragraph_ledger.json final_prose. Do NOT hand-edit.

\chapter{Main Empirical Analyses}\label{ch:main}

\section{Data, Sample, and Variable Construction}

The analysis combines five data sources. Capital IQ supplies the earnings-call transcripts of United States public firms over 2002 to 2018, parsed to speaker role and to the presentation and question-and-answer segments of each call. Quarterly Compustat North America fundamentals supply firm financials, taken from the most recent fiscal quarter ending on or before the call. CRSP daily records and IBES earnings records supply the stock-return, market-return, and earnings-surprise controls that the decomposition of Section~2.3 requires. SDC supplies the merger and acquisition records, including each deal's announcement date, its completion or withdrawal date, and its cash-versus-stock payment composition. The identity of the chief executive on each call, the input to the manager fixed effect of Section~2.3, comes from Execucomp, whose annual records we aggregate into a monthly panel of which executive holds the CEO role in each quarter. These sources are keyed to different identifier systems, so they must be linked: SDC keys deals by the acquirer's six-digit CUSIP while the call panel is keyed by Compustat firm identifier (gvkey), and we join deals to calls through the acquirer's six-digit CUSIP via the CRSP-Compustat link table recorded in the master-sample manifest. The cash-versus-stock classification and the event clock that the designs below depend on exist only through this match, which we record as a researcher degree of freedom rather than assume away.

The uncertainty measures are built in-house from the language of the call. Each transcript is tokenized, and the Loughran and McDonald uncertainty-word share is computed separately by speaker and by segment. The CEO's question-and-answer share, $\mathrm{UncAnsCEO}$, is the input that the first stage of the Section~2.3 decomposition, Equation~(4) of \citet{dwz} estimated on our sample, residualizes into $\mathrm{UncResCEO}$, the call-varying component net of the CEO's persistent speaking style. What we discharge here is the construction mechanic only: the equation itself, and the definition and justification of the residual, are set out in Section~2.3 and are not re-derived.

The remaining key estimation variables are constructed and pointed back to their defining homes rather than re-derived here. $\mathrm{UncResCEO}$ is defined in Section~2.3. $\mathrm{PreAnnounceQtr}=\mathbf{1}[e=-1]$ marks the quarter before a firm's first qualifying deal, as defined in Sections~2.2 and~2.4. $\mathrm{CashRatio}$ is cash and equivalents over total assets, $\mathrm{cheq}/\mathrm{atq}$, as defined in Section~2.5 and the Appendix. $\mathrm{CashScrutiny}$ is the share of analyst question-and-answer turns devoted to cash and liquidity, read from the locked word list in Appendix~I, and $\mathrm{HighCashScrutiny}$ flags any such scrutiny at all, since the median call draws none; both are defined in Section~2.5, and only about eleven percent of calls draw cash scrutiny (the $\mathrm{HighCashScrutiny}$ mean is $0.1127$, Panel~A). The firm-financial controls, leverage, log assets, Tobin's $Q$, return on assets, capital expenditure, a dividend indicator, and cash-flow volatility, are catalogued in the Appendix.

The estimation samples are built in three layers, and it is this layering that makes the number of observations vary from one table to the next. The base call panel comprises every matched Capital IQ call of a United States public, non-financial, non-utility firm over 2002 to 2018, excluding CEOs with fewer than five calls, for whom no speaking style and hence no residual can be estimated. Because the manager fixed effect of Section~2.3 identifies each CEO from Execucomp, the sample is bounded to the firms it covers, approximately the S\&P~1500, which skews the uncertainty regressions toward larger, more heavily covered firms. On top of it sits the SDC deal universe: United States public acquirers over the same period with deal status in completed, pending, or withdrawn and a known payment composition, where deals at least half in cash form the cash arm and deals at least half in stock the stock placebo arm, each firm's first qualifying deal setting its event clock. The third layer is a set of design-specific trims: the run-up test drops treated firms' post-announcement quarters, the event studies cap the post-window at four quarters and truncate it at the firm's next announcement, and the scrutiny validity and channel tests require at least three analyst question-and-answer turns per call. Each test thus keeps only the firms for which its particular comparison is defined. The five-call requirement enters here as a build mechanic; its consequences for what the residual subsample represents are treated as a selection question in Section~2.4.

Table~\ref{tab:summary_stats} reports summary statistics for the main estimation sample, an anchor of $88{,}205$ earnings-call observations across $1{,}884$ firms over 2002 to 2018. The number of observations varies by row because each variable is summarized on the complete cases of its own anchor panel. The residual $\mathrm{UncResCEO}$ is centered near zero, with a mean of $-0.0000$ and a standard deviation of $0.3010$ (Panel~B); the cash ratio has a mean of $0.1708$ (Panel~C) and the pre-announcement indicator a mean of $0.0094$ (Panel~A). Because the residual exists only on the at-least-five-call, non-financial and non-utility universe, the uncertainty regressions describe a more-covered set of firms, skewed toward larger ones, than the cash-ratio regressions, and for that reason every table that follows reports its own firm and observation counts. We read these statistics descriptively, as the foundation the results sections build on rather than as findings in themselves.

\section{Main Analysis 1: The Pre-Announcement Run-Up}

The first main analysis tests H1 with the single-indicator, two-way fixed-effects regression of Section~2.4, carrying firm and calendar year-quarter fixed effects and standard errors clustered by firm. It is estimated four times, once for each outcome, the cash ratio and the residual $\mathrm{UncResCEO}$, in each arm, the cash acquirers (at least half in cash) and the stock acquirers (at least half in stock) that serve as a placebo. Treated firms' post-announcement quarters are dropped, so the design observes only the run-up and not the aftermath. Identification is within firm, setting a treated firm's pre-announcement quarter against its own ordinary quarters; the estimating equation is the one given in Section~2.4 and is not restated. The coefficient on $\mathrm{PreAnnounceQtr}$ estimates the pre-announcement shift $\theta_{-1}$, and we read it descriptively, as an anticipation contrast, with no causal or parallel-trends claim.

In the cash arm both clocks move in the final pre-announcement quarter. Residual question-and-answer uncertainty is elevated (coefficient $0.0461$, $p<0.01$; column~2 of Table~\ref{tab:empire_building_did}), and the cash ratio is rising (coefficient $0.0051$, $p<0.01$; column~1). The cash-ratio coefficient tracks the change in cash rather than its level, because that regression carries the cash ratio's own one-quarter lag as a partial-adjustment term (coefficient $0.7990$, $p<0.01$; column~1). The residual elevation is the load-bearing result, and although the table reports the treatment coefficient one-tailed, it survives a two-tailed test as well ($t=2.68$, $p=0.0074$), so the result does not rest on the reporting convention.

The two cash-arm columns are not estimated on the same set of firms, because the residual exists only for CEOs with at least five calls. The uncertainty test in column~2 uses $27{,}622$ firm-quarters across $1{,}248$ firms, whereas the cash-ratio test in column~1 uses $67{,}590$ firm-quarters across $2{,}232$ firms. The uncertainty test is therefore the more demanding of the two, run on the smaller and more heavily covered universe, a comparability limit we disclose so the two columns are not read as a single sample.

The effect is material but modest. The uncertainty coefficient equals roughly fifteen percent of a residual standard deviation, about $15.3$ percent against the standard deviation of $0.3010$ in Panel~B of Table~\ref{tab:summary_stats}, for a variable from which everything persistent has already been removed, and it appears in a quarter that is itself about one percent of the sample (the mean of $\mathrm{PreAnnounceQtr}$ is $0.0094$, Panel~A). The cash coefficient adds roughly half a percentage point of assets to the cash stock in a single quarter, about $3.0$ percent of the mean cash ratio of $0.1708$ (Panel~C). The magnitudes are economically meaningful without being large, which is the register in which we read them.

The stock arm is the placebo. Under the same disclosure constraint, and with an otherwise nearly identical model, neither outcome shows a comparable positive run-up: residual uncertainty is, if anything, lower (coefficient $-0.0429$, not significant; column~6 of Table~\ref{tab:empire_building_did}) and the cash ratio is flat (coefficient $-0.0036$, not significant; column~5). The two models are close to identical in their dynamics, the stock partial-adjustment lag of $0.8013$ standing almost exactly where the cash lag of $0.7990$ does, so what differs at the pre-announcement quarter is the response and not the specification. A side-by-side reading of this kind is not yet a formal test of the cash-versus-stock difference; that test is supplied in Section~3.4, the third main analysis.

\section{Main Analysis 2: Differential Timing Around the Announcement}

The second main analysis replaces the single pre-announcement indicator with the four-bin event study of Section~2.4. The bins are $\mathrm{PRE2}$ at two quarters before announcement, $\mathrm{PRE1}$ at one quarter before, $\mathrm{GAP}$ for the interval in which the deal is announced but not yet closed, and $\mathrm{POST}$ for the quarters after completion, with the baseline formed by quarters three or more before announcement together with never-acquirers. The specification is otherwise that of the first main analysis, and because no ordering is imposed across the bins every coefficient is reported two-tailed. The test is run on the matched universe in which both the residual and the cash ratio are defined, $28{,}102$ firm-quarters across $1{,}320$ firms, so that both outcomes are estimated on the identical rows; estimating them on the same rows is what lets the two paths be read against each other while guarding the differential-timing reading against a sample-composition artifact. The estimating equation is the one given in Section~2.4 and is not restated; the reading remains within firm and correlational, with no identification claim.

The pre-trend is flat on both clocks. Two quarters before announcement, at $\mathrm{PRE2}$, the coefficient is null for the residual (coefficient $0.0068$, not significant; column~1 of Table~\ref{tab:empire_drop_matched}) and for the cash ratio (coefficient $0.0008$, not significant; column~2), so the one-quarter-before estimates are not the tail of a pre-existing trend. This is a validity device, not a proof of identification.

The discriminating contrast is the drop at the gap. On the information clock, residual uncertainty spikes one quarter before announcement (coefficient $0.0473$, $p<0.01$; column~1 of Table~\ref{tab:empire_drop_matched}) and falls to insignificance the instant the deal is public (coefficient at $\mathrm{GAP}$ of $0.0018$, not significant), a significant decline from the pre-announcement peak to the gap (drop $0.0455$, $p<0.05$) and a large significant swing from the peak to completion (drop $0.0723$, $p<0.01$). On the transaction clock, the cash ratio is elevated one quarter before announcement (coefficient $0.0061$, $p<0.05$; column~2), does not fall when the deal is announced (the pre-announcement-to-gap drop of $0.0006$ is not significant), and declines only once the cash leaves at completion (drop from gap to completion $0.0210$, $p<0.01$; the completion coefficient itself is $-0.0155$, $p<0.01$). Because the gap-quarter cash coefficient itself, $0.0055$, is positive but not significant, the persistence of cash rests on the absence of a decline at announcement rather than on a significant gap-quarter level. The contrast between these two paths is the claim, read within firm and correlationally.

The pre-announcement spike is economically the same event as the one the first main analysis measured: $0.0473$ here against $0.0461$ there, about fifteen percent of a residual standard deviation, $15.7$ percent against the $0.3010$ in Panel~B of Table~\ref{tab:summary_stats}. What the event study adds is the round trip. The swing from the pre-announcement peak to completion, a drop of $0.0723$ ($p<0.01$), shows that the uncertainty raised in the pre-announcement window is fully unwound once the information is public. The cash leg, for its part, is identified off the same well-behaved partial-adjustment specification used in the first main analysis, its one-quarter lag standing at $0.7547$ ($p<0.01$; column~2 of Table~\ref{tab:empire_drop_matched}). Because this event study runs on the matched universe rather than the all-firm sample, scaling by the all-universe standard deviation is a disclosed approximation, and the magnitude is read as supportive rather than decisive.

The same round trip appears in the cash arm of the placebo event study, estimated on its own complete-case cash universe of $29{,}535$ firm-quarters across $1{,}326$ firms. There residual uncertainty again spikes one quarter before announcement (coefficient $0.0486$, $p<0.01$; column~1 of Table~\ref{tab:empire_drop_placebo}) and collapses to insignificance at the gap (coefficient $0.0058$, not significant), with a significant decline from the peak to the gap (drop $0.0428$, $p<0.10$) and a significant swing from the peak to completion (drop $0.0681$, $p<0.01$). This is an independent corroboration of the matched-universe result on a second sample. Only the cash column of this placebo table is used here; its cash-versus-stock contrast is reserved for Section~3.4.

The claim rests on the contrast between the two paths, not on the magnitude or the sign of any single drop. Uncertainty tracks the information and reverses at announcement; cash tracks the transaction and falls at completion, a persistence that is in part mechanical, since the cash leaves the balance sheet only at closing. We do not over-read the below-baseline completion coefficient on the uncertainty clock (coefficient $-0.0250$, $p<0.10$; column~1 of Table~\ref{tab:empire_drop_matched}); it is flagged rather than interpreted. The reading throughout is correlational and descriptive, supportive of the disclosure-state pattern rather than definitive about any mechanism, which we leave open.

\section{Main Analysis 3: Cash-Specificity}

A first step reruns the event study side by side on the stock placebo arm, using the stock column of Table~\ref{tab:empire_drop_placebo} that was held back from the previous analysis. Neither arm trends into the event: the pre-trend coefficient two quarters out is small and insignificant in both, $0.0105$ for the cash arm and $-0.0056$ for the stock arm. From there the two diverge. The cash arm reproduces the run-up, with a spike one quarter before announcement (coefficient $0.0486$, $p<0.01$) and a swing from the peak to completion (drop $0.0681$, $p<0.01$), while the stock arm shows none of it (coefficient one quarter before announcement of $-0.0404$, not significant) and, if anything, moves the other way across announcement (decline from the pre-announcement quarter to the gap of $-0.0756$, $p<0.10$). This is a descriptive comparison; it motivates a formal test rather than standing in for one.

A side-by-side display of this kind implies a comparison it never actually runs. The formal test instead uses the pooled, interacted model of Section~2.4, which carries both treatments at once and tests the cash-minus-stock difference directly through a single Wald linear restriction, $\beta_c-\beta_s>0$. That restriction is what estimates $\theta_{-1}^{\,\mathrm{cash}}-\theta_{-1}^{\,\mathrm{stock}}$ in H1a, and it is this difference, not a comparison of one significant coefficient against one insignificant one, that constitutes the cash-specificity test. The pooled baseline is the set of firms that make neither a cash nor a stock acquisition. The reading remains correlational, with no identification claim, and the framing is one of concentration rather than strict specificity.

On the effect, the pooled model delivers the formal cash-specificity result. The cash arm's pre-announcement uncertainty coefficient is significant (coefficient $0.0459$, $p<0.05$; column~1 of Table~\ref{tab:empire_cashspec}) and the stock arm's is insignificant and negative (coefficient $-0.0524$, not significant). The formal cash-minus-stock difference, evaluated as the single Wald restriction, is $0.0983$ (standard error $0.0476$, $z\approx2.07$, $p=0.039$), significant at the five percent level on a two-tailed test. Because the two arms move in opposite directions, the difference is larger than either coefficient, roughly a third of a residual standard deviation, about $32.7$ percent against the $0.3010$ in Panel~B of Table~\ref{tab:summary_stats}. It is this formal linear restriction, and not the contrast between a significant and an insignificant coefficient, that carries the cash-specificity claim.

On the proposed cause, the picture is weaker. The cash build-up itself is only faintly cash-sided: the cash-minus-stock difference in the cash ratio is insignificant on the matched universe (difference $0.0064$, not significant; column~2 of Table~\ref{tab:empire_cashspec}) and reaches only marginal significance on the full panel (difference $0.0092$, $p<0.10$; column~3). The effect concentrates in cash deals, but the war-chest mechanism that might explain it is therefore not established, and the mechanism is left open.

We therefore read the result as supported but fragile. One formal test at the five percent level supports H1a's claim of concentration in cash, but its significance rides on the imprecise, negative stock estimate (coefficient $-0.0524$, standard error $0.0436$; column~1 of Table~\ref{tab:empire_cashspec}), and the cause's own cash-specificity is weaker still, marginal and only on the full panel. The wording stays with concentration rather than strict specificity, the difference between a significant and an insignificant coefficient is not treated as the test, and no cause is asserted.

\chapter{Additional Analyses}\label{ch:additional}

\section{Ruling Out Analyst Scrutiny}

The most immediate rival to the anticipatory reading is analyst scrutiny: that the quarters before a deal draw harder analyst questions about cash, and that the hard questioning, rather than the deal itself, produces the hedged answers. This is a genuine confound and not a strawman. The incidence of cash scrutiny itself rises in the quarter before a cash acquisition (the indicator $\mathrm{HighCashScrutiny}$ carries a coefficient of $0.0408$, $p<0.05$; column~4 of Table~\ref{tab:empire_building_did}), so it co-moves with both the cash position and the event and could in principle generate the run-up mechanically. The continuous volume of cash questions rises too but not significantly (coefficient $0.0854$, not significant; column~3), so what rises reliably before a cash deal is the incidence of cash scrutiny, not its volume. We address this rival in three steps.

The first step is to confirm that the scrutiny measure is real. The share of analyst question-and-answer turns devoted to cash and liquidity rises strongly with the firm's cash ratio (coefficient $0.7530$, $p<0.01$, in a one-tailed test, rising to $0.8519$, $p<0.01$, with the full control set; columns~1 and~2 of Table~\ref{tab:cash_scrutiny_validity}) and with a high-cash indicator (coefficient $0.1754$, $p<0.01$; column~3), so it tracks what it is built to track and the nulls that follow are not the artifact of a dead measure.

The second step asks whether scrutiny alone moves the residual, and it does not. Its direct effect on $\mathrm{UncResCEO}$ is zero to four decimal places (coefficient $-0.0000$, not significant; standard error $0.0013$; $41{,}512$ calls; Panel~A of Table~\ref{tab:cash_scrutiny_channel}), and the null holds under a logit on the dichotomized residual (log-odds coefficient $-0.0003$, not significant; Panel~B) and under the raw uncertainty measure and coarser scrutiny proxies alike. This test is not redundant with the way the residual is built: $\mathrm{CashScrutiny}$ is the volume of cash questions, whereas the uncertainty wording of those questions, $\mathrm{UncQue}$, is already netted out inside the residual by the decomposition of Section~2.3, so the null speaks to a path the residual does not already absorb.

The third and operative step puts scrutiny and the pre-announcement indicator in the same model, estimated over the cash acquirers' pre-announcement and earlier quarters. The pre-announcement indicator stays significant, both with scrutiny entered (coefficient $0.0413$, $p<0.01$; column~1 of Table~\ref{tab:reason_gating}) and with the scrutiny-by-pre-announcement interaction added (coefficient $0.0439$, $p<0.05$; column~2), while scrutiny's own effect remains null ($0.0000$ and $0.0001$, neither significant) and the interaction is insignificant (coefficient $-0.0056$, not significant; column~2). The reason raises residual uncertainty; the questioning does not, and the questioning does not amplify the reason.

The register here should stay honest. This is a failure to find, paired with a within-model dissociation, and not a powered test of equivalence. The interaction's confidence interval, approximately $[-0.027, 0.016]$, does not exclude effects that would matter against a run-up near $0.04$; about eighty-nine percent of calls draw no cash scrutiny at all; and the contrast between a significant and an insignificant coefficient is not itself a test. The claim is that scrutiny does not account for this particular run-up, not that scrutiny never matters anywhere.

